\section{方法}
\label{sec:method}

\subsection{总体结构}

CogECI 的整体流程如图~\ref{fig:frame} 所示。输入是一篇文档的句子序列与其中的事件提及；
输出是每个候选事件对的因果标签。整条链路包含四个部分：句子编码与句子图模块（SGM）、
事件表示与双流因果推理（DSCRM）、迭代式因果图更新（继承自 iLIF）、
以及上下文对齐损失（CGL）。

\begin{figure}[H]
\centering
\begin{tikzpicture}[
  x=1cm, y=1cm,
  box/.style={draw=deepblue!70, fill=deepblue!6, rounded corners=2pt,
              minimum height=0.86cm, text width=5.2cm, align=center,
              font=\small, inner sep=3pt},
  loss/.style={draw=orange!75!black, fill=orange!10, rounded corners=2pt,
               minimum height=0.86cm, text width=3.9cm, align=center,
               font=\small, inner sep=3pt},
  arr/.style={-{Stealth[length=2mm]}, draw=black!70, thick},
  darr/.style={-{Stealth[length=2mm]}, draw=orange!75!black, thick, dashed},
  larr/.style={-{Stealth[length=2mm]}, draw=deepblue!70, thick, dashed}
]
% ---- 主干：自上而下 ----
\node[box] (doc)  at (0, 0)    {文档句子 $s_1,\dots,s_n$};
\node[box] (enc)  at (0,-1.30) {BERT 编码 + SGM 句子图 GCN};
\node[box] (evt)  at (0,-2.60) {事件提及 $\rightarrow$ 片段内子词均值 $\rightarrow$ 事件表示 $\bm{e}_k$};
\node[box] (pair) at (0,-3.90) {事件对表示 $[\bm{e}_i\,\|\,\bm{e}_j\,\|\,(\bm{z}_i-\bm{z}_j)]$};
\node[box] (mlp)  at (0,-5.20) {DSCRM 双流 MLP（句内 / 跨句参数独立）};
\node[box] (pred) at (0,-6.50) {因果预测 $\hat{\bm{p}}_{ij}$};
\node[box] (cge)  at (0,-7.80) {置信度 $>\omega$ 建因果图 $\rightarrow$ CGE 更新 $\bm{z}$};

\foreach \a/\b in {doc/enc, enc/evt, evt/pair, pair/mlp, mlp/pred, pred/cge}
  \draw[arr] (\a) -- (\b);

% ---- 右侧：迭代回路 ----
\coordinate (r1) at (4.35,-7.80);
\coordinate (r2) at (4.35,-3.90);
\draw[larr] (cge.east) -- (r1) -- (r2) -- (pair.east);
\node[align=left, font=\scriptsize, text=deepblue!75, anchor=west] at (4.5,-5.85)
  {迭代更新 $\bm{z}$\\（至 $L_{\max}$ 轮，\\结构差 $\le\delta$ 早停）};

% ---- 左侧：两个损失 ----
\node[loss] (ecl) at (-4.75,-3.90) {ECL：Focal Loss\\$\gamma=2,\ \alpha=0.75$};
\node[loss] (cgl) at (-4.75,-5.55) {CGL：句级 BCE\\$w=6$};
\draw[darr] (pair.west) -- (ecl.east);
\draw[darr] (pair.west) -- ++(-0.55,0) |- (cgl.east);

\end{tikzpicture}
\caption{CogECI 的结构示意。黑色实线为前向数据流；橙色虚线为两个损失项，
  以事件对表示为监督对象；蓝色虚线为迭代回路，每轮用置信度 $>\omega$ 的预测重建因果图，
  经 CGE 更新事件表示 $\bm{z}$ 后回到事件对表示。}
\label{fig:frame}
\end{figure}
